\subsection{A permitted auxiliary prime with an incompatible norm unit}

\begin{statement}{A.103}{Proved}
{The coupled residue equations survive boundary and prior-radical splitting}
Let $K=\Q(\alpha)$, $\alpha^3=3$, and put
\[
 \begin{gathered}
 q=22441,\qquad c=4048,\qquad H=c^3-3=q\cdot2955829,\\
 \beta_0=1-\alpha,\qquad
 u=\frac{(c-\alpha)^3}{H},\qquad \beta=\beta_0u.
 \end{gathered}
 \tag{A.103.1}
\]
Then $N_{K/\Q}u=1$ and $N_{K/\Q}\beta=-2$.
The prime $q$ splits completely in
\[
 \Q\bigl(\mu_{24},(-2)^{1/24},\alpha,
                 (1-\zeta_3^i\alpha)^{1/3}:0\le i<3\bigr),
 \qquad v_3(q-1)=1.
 \tag{A.103.2}
\]
Thus it satisfies both the full boundary condition at
exponent $24$ and the conjugate cubic-radical splitting
condition for the \emph{prior} representative $\beta_0$.
The change from $\beta_0$ to $\beta$ preserves its canonical
two-adic divisor from A.101.

For every integer $e$ prime to three, every
$Q\in\Q^\times$, and every $R\in\Q^\times$ with
$v_q(R)=0$, the actual cubic extension
\[
 E_Q=K(\gamma),\qquad \gamma^3=Q^e\alpha
 \tag{A.103.3}
\]
has the following property:
$\beta R^e$ is not a norm from $E_Q^\times$ to $K^\times$.
It is not a norm between their three-primary completions
either.
There is no restriction on $k=v_q(Q)$; in particular
allocating $q$ to the updated $A$-coefficient does not
repair this specified representative.
For the coefficient update with $M=24$ one takes $e=8$,
so $a'=3Q^{24}$ and $b'=-2R^{24}$.
\end{statement}
\begin{proof}
The three roots of $T^3-3$ in $\F_q$ are
\[
       a_0=4048,\qquad a_1=8110,\qquad a_2=10283.
 \tag{A.103.4}
\]
Direct modular exponentiation gives
\[
 \begin{aligned}
 q&\equiv49\pmod{72},&
 (-2)^{(q-1)/24}&=1,\\
 (1-a_i)^{(q-1)/3}&=1\quad(0\le i<3),&
 \bigl((a_i/a_0)^{(q-1)/3}\bigr)_{i=0}^2
     &=(1,20414,2026).
 \end{aligned}
 \tag{A.103.5}
\]
The prime and the complete list of roots, as well as these
finite identities, are checked by the retained exact
certificate below.
They imply (A.103.2): all the displayed defining
polynomials split with distinct nonzero roots in
$\F_q$, and $q$ is away from their discriminant primes.
The last two values in (A.103.5) are the two nontrivial
cube roots of unity in $\F_q$.

The norm identity $N(c-\alpha)=c^3-3=H$ gives the
claimed norms.
For a literal polynomial representative,
\[
 \beta=\frac{
  66331722157-66380917501\alpha+49171056\alpha^2}
 {66331758589}.
 \tag{A.103.6}
\]
At two, $c$ is even and $H$ is odd.
Each conjugate of $\alpha$ is a two-adic unit, so $u$
is a unit at every prime over two.
The two-adic divisor of $\beta_0$ is therefore retained:
the unique linear root is $3$ modulo eight, giving
$v_2(1-\alpha)=1$ there, and the other roots reduce
to the nontrivial cube roots of unity in $\F_4$,
where $1-\alpha$ is a unit.

Let $\alpha_i\in\Q_q$ denote the Hensel lifts of
(A.103.4).
Since $v_q(H)=1$, one has
$v_q(c-\alpha_0)=1$ and
$v_q(c-\alpha_1)=v_q(c-\alpha_2)=0$.
The three factors $1-\alpha_i$ are units and are
cubes in $\Q_q^\times$ by (A.103.5) and Hensel lifting.
Consequently the valuation vector and the cubic
classes of the unit parts of $\beta$ are
\[
 (t_0,t_1,t_2)=(2,-1,-1),\qquad
 [q^{-t_i}\beta_i]=[(H/q)^{-1}]
       \quad\hbox{in }\F_q^\times/(\F_q^\times)^3.
 \tag{A.103.7}
\]
In particular the three unit classes coincide.
Their common cubic character is $2026$.

We specify the local norm criterion used here.
For $F=\Q_q$ with $\mu_3\subset F$, write
$d=q^s d_0$ and $z=q^t z_0$, with $d_0,z_0$ units.
If a cubic root of $d$ is adjoined, the norm condition is
\[
       t[d_0]-s[z_0]=0
          \quad\hbox{in }\F_q^\times/(\F_q^\times)^3.
 \tag{A.103.8}
\]
This includes the split cubic algebra.
Here and below the residue multiplicative group is
written additively.
For completeness, if $3\mid s$, reduce $d$ by a cube.
A noncube unit defines the unramified cubic extension;
its units are norms and its norm valuations are
multiples of three, giving exactly (A.103.8).
If $3\nmid s$, choose integers $a,b$ with $as-3b=1$.
The element $(d^{1/3})^a/q^b$ is a uniformizer in
the totally ramified cubic extension, and its norm is
$q d_0^a$.
The residues of norms of units are precisely cubes:
necessity follows by multiplying the three equal
residues, and sufficiency follows because the cube
map is bijective on the principal units of $F$.
Thus the norm condition is
$[z_0]=at[d_0]$, equivalent to (A.103.8).
This proves the formula without suppressing either
the valuation or its associated unit.

Write $Q=q^kQ_0$, and set
\[
 x=[Q_0],\quad y=[R],\quad A_i=[\alpha_i],
       \quad C=[(H/q)^{-1}]
             \quad\hbox{in }\F_q^\times/(\F_q^\times)^3.
 \tag{A.103.9}
\]
At branch $i$, the radicand in (A.103.3) has valuation
$ek$ and unit class $ex+A_i$.
The target $\beta_iR^e$ has valuation $t_i$ and
unit class $C+ey$.
Equation (A.103.8) therefore requires the three exact
linear equations
\[
             t_i(ex+A_i)-ek(C+ey)=0,
                           \qquad 0\le i<3.
 \tag{A.103.10}
\]
All $t_i$ reduce to $-1$ modulo three.
Subtracting two equations removes $e,k,x,y,C$ and
forces $A_i=A_j$.
This contradicts the distinct root-ratio classes in
(A.103.5).
The contradiction applies to every integer $k$,
including zero and multiples of three.

The extension (A.103.3) really has degree three:
at the unique prime of $K$ above three, the valuation
of its radicand is $1+3e\,v_3(Q)$, which is not
divisible by three.
A global norm must satisfy the local norm condition
at every split $q$-branch, so the failure just proved
excludes a global norm.
The failed local norm quotient has exponent three
and is a finite continuous quotient of $\Q_q^\times$.
It therefore also factors through its three-primary
completion and excludes a completed norm.
\end{proof}

The exact updated norm-fibre map is
\[
 \left\{z\in\widehat{E_Q^\times}_3:
       N_{E_Q/\Q}z=-2R^{3e}\right\}
       \xrightarrow{\ N_{E_Q/K}\ }
 \left\{w\in\widehat{K^\times}_3:
       N_{K/\Q}w=-2R^{3e}\right\}.
 \tag{A.103.11}
\]
Its target contains $\beta R^e$, which is outside
its image.
This is a restriction on the specified representative
under this update, not a nonvanishing assertion for
the underlying pair: for $Q=R=1$, the elements
$1-3^{1/3^j}$ give the exact compatible norm family
with rational norm $-2$.
In particular the splitting conditions in A.101 remain
available, but they do not specify the post-repair
unit classes needed by (A.103.10).

\noindent The executed certificate
\path{code/norm_permitted_unit_certificate.py}, with output
\path{code/norm_permitted_unit_certificate_output.txt}, checks
primality, all splitting residues, the exact polynomial
norms, the local valuation and unit characters, and all
$54$ choices of $e,k,x,y$ modulo three with $3\nmid e$.
It also checks directly that the modification is a unit
at every place above two.
